\section{Multi-Agent System Architecture}\label{sec:architecture}
The global pipeline is organized around three specialized agents coordinated by a central orchestrator. Figure~\ref{fig:pipeline_professional_colored} summarizes the end-to-end data flow and feedback loops in the iterative process.

\begin{figure*}[t]
\centering
\begin{tikzpicture}[
	>=latex,
	% Base Agent Style
	agent/.style={
		rectangle,
		thick,
		minimum width=3.5cm,
		minimum height=1.2cm,
		align=center,
		font=\sffamily\bfseries,
		rounded corners=2pt
	},
	% Distinct Professional Colors for each Agent
	orchStyle/.style={agent, fill=purple!10, draw=purple!80!black, minimum width=5cm},
	genStyle/.style={agent, fill=orange!10, draw=orange!80!black},
	chemStyle/.style={agent, fill=blue!10, draw=blue!80!black},
	bioStyle/.style={agent, fill=green!10, draw=green!80!black},
	% Input / Output Styles
	io/.style={rectangle, thick, align=center, font=\sffamily\footnotesize, rounded corners=2pt},
	inputStyle/.style={io, dashed, fill=gray!5, draw=gray!60},
	outputStyle/.style={io, solid, fill=yellow!15, draw=yellow!60!black},
	% Flow Lines
	batchFlow/.style={-{Stealth[length=3mm]}, thick, draw=black!90},
	scoreFlow/.style={-{Stealth[length=3mm]}, thick, dashed, draw=black!70}
]

% --- NODES ---
\node[orchStyle] (orch) at (0, 3) {Orchestrator};
\node[genStyle]  (gen)  at (-5, 0) {Generator Agent};
\node[chemStyle] (chem) at (0, -3) {Chemistry Agent};
\node[bioStyle]  (bio)  at (5, 0) {Biology Agent};

% --- SYSTEM INPUTS & OUTPUTS ---
\node[inputStyle] (inputs) at (-2.5, 5) {\textbf{Inputs}: Hyperparameters\\(Top-$K$, Iterations, Batch Size, Targets)};
\node[outputStyle] (outputs) at (2.5, 5) {\textbf{Output}: Final\\Top-$K$ Peptides};

% Connect IO to Orchestrator (Straight up/down)
\draw[batchFlow] (inputs) -- (inputs |- orch.north);
\draw[batchFlow] (outputs |- orch.north) -- (outputs);

% --- BATCH DATA FLOW (Solid Lines) ---
% 1. Orchestrator -> Generator
\draw[batchFlow] (orch.west) -| node[pos=0.25, above, font=\sffamily\footnotesize, align=center] {Init / Target Stats\\\& Ref. Peptide} (gen.north);

% 2. Generator -> Chemistry
\draw[batchFlow] (gen.south) |- node[pos=0.75, above, font=\sffamily\footnotesize, align=center] {Candidate Batch} (chem.west);

% 3. Chemistry -> Biology
\draw[batchFlow] (chem.east) -| node[pos=0.25, above, font=\sffamily\footnotesize, align=center] {Filtered Batch} (bio.south);

% --- FEEDBACK & SCORE FLOW (Dashed Lines) ---
% 4. Chemistry -> Orchestrator (Straight up through the center)
\draw[scoreFlow] (chem.north) -- node[fill=white, inner sep=2pt, font=\sffamily\footnotesize, align=center] {Chem Scores\\(Validity \& Proxy Stats)} (orch.south);

% 5. Biology -> Orchestrator
\draw[scoreFlow] (bio.north) |- node[pos=0.75, above, font=\sffamily\footnotesize, align=center] {Bio Scores\\(Activity \& Risks)} (orch.east);

% Enforce a symmetric final bounding box around x=0 to guarantee true global centering.
\pgfresetboundingbox
\path (-5.9, -3.4) rectangle (8.2, 3);

\end{tikzpicture}
\caption{Agentic Pipeline Architecture. Solid lines denote the sequential flow of peptide candidate batches. Dashed lines indicate the feedback loops, where both the Chemistry and Biology agents return their respective proxy scores back to the Orchestrator. Color coding differentiates the specialized domains of the pipeline.}
\label{fig:pipeline_professional_colored}
\end{figure*}

\subsection{The Orchestrator}\label{subsec:orchestrator}
The Orchestrator is the control layer that turns the three experts into a single iterative optimization system. 
Its role is to enforce a fixed sequence of operations:
\begin{enumerate}
	\item Request a batch of candidates from the Generator (if not given when originally called)
	\item Pass this batch through the Chemistry Agent to keep only chemically admissible peptides and associated proxy descriptors, and  send the surviving candidates to the Biology Agent for biological scoring. 
	\item Based on these expert outputs, the Orchestrator ranks candidates, keeps the Top-$K$ subset, and uses this subset as the seed for the next generation cycle.
\end{enumerate}
In this way, each iteration combines exploration (new proposals) and exploitation (retaining best current candidates) under one consistent decision policy.
The stopping criterias are handled at the orchestration level rather than inside each expert module. \\In practice, the loop can stop when :
\begin{itemize}
	\item A maximum number of iterations is reached
	\item No valid candidates survive the chemistry gate
	\item The optimization objective is considered stable enough according to user-configured run settings (for example fixed Top-$K$, batch size, and iteration budget)
	%TODO UPDATE APRES LA PR DE L'ORCHESTRATOR
\end{itemize}
This implementation allows reproducible exepimentations : changing run behavior only requires updating orchestrator-level parameters, there is no need to change the experts.
This agent also provides traceability by recording round-wise pipeline events and outcomes:
\begin{itemize}
	\item Input settings
	\item Filtered populations
	\item Expert scores
	\item Selected candidates
	%TODO UPDATE APRES LA PR DE L'ORCHESTRATOR
\end{itemize}
This log-centric design makes the optimization path auditable from end to end: one can reconstruct why a peptide was retained or discarded at a given step, compare runs under different settings, and support post-hoc analysis of convergence behavior and failure cases.

\subsection{Generator Agent}\label{subsec:designer}
TODO: Describe how this agent generates candidates (prompting strategies, templates, mutation/crossover operators)[cite: 15].

\subsection{Chemistry Agent}\label{subsec:chemistry}
The Chemistry Agent is the deterministic constraint-and-scoring stage of the pipeline. At a global level, it receives peptide sequences from the generator, evaluates them against user-defined physicochemical targets, and returns both a feasibility decision and a quantitative score used downstream by the orchestrator. The behavior of this agent is driven by a configurable target profile (ranges, target values, optional weights, and pH), which allows the same workflow to be reused under different experimental objectives.

Its workflow can be summarized through three complementary functions:

\noindent\textbf{1. Sanitization :}
Generative models can hallucinate sequences that are syntactically correct but chemically invalid or unstable. The Chemistry Agent goes beyond simple string matching by attempting to instantiate a molecular graph for each sequence using the RDKit framework.
This process involves a full sanitization step: checking valence rules, aromaticity, and atom consistency. Any sequence that fails to convert into a valid RDKit Mol object is immediately discarded as a chemical syntax error. This ensures that downstream agents only process realizable matter.

\medskip
\noindent\textbf{2. Multidimensional Property Profiling :}
To characterize developability, the agent computes a configurable set of descriptors from two complementary sources: peptide-centric properties (length, net charge at chosen pH, isoelectric point, hydrophobicity, cationicity) and molecule-centric properties (notably RDKit-based \(\log P\)). This mixed profile lets the orchestrator compare candidates in a unified physicochemical space.

\medskip
\noindent\textbf{3. Quantitative Scoring and Adaptive Filtering :}
The agent assigns each candidate a continuous score based on distance to target values, with optional per-property weighting. It then applies a two-step filtering policy: first keep peptides that satisfy all hard limits, then (if too few survive) refill the population with near-miss candidates that are closest to the target profile. This adaptive mechanism prevents early pipeline collapse while preserving optimization pressure.
\\

The Chemist agent is the first called agent after the generator is used. It allows a first selection and analysis on the generated peptides list before sending them to the Biologist agent.
Based on the strict limits provided by the user it sends back to the orchestrator scores and properties on the peptides. From an higher-level perspective, the returned data acts as a deterministic generated ground-truth wich can be trusted by the Biologist agent for its prediction.

\subsection{Biology Agent (Function \& Risk)}\label{subsec:biology}
The Biology Agent evaluates the generated candidates and provides the biological signal used by the iterative loop. In our implementation, it is based on Meta/Facebook's Evolutionary Scale Modeling family (ESM-2) \cite{rives2021biological,esmgithub} and can be conditioned by context (for example an existing peptide reference or requested target properties).

At each iteration, the Biology Agent scores the current valid batch and returns values used by the orchestrator to rank candidates and keep the Top-$K$ peptides for the next round. This makes it the core selection component on the biological side of the pipeline.

As a design choice, this score is treated as a relative proxy to guide search, not as a direct experimental measurement. The detailed scoring formulation is reported in Section~\ref{subsec:scoring}.
